%%%%%%%%%%%%%%%%%%%%%%%%%%%%%%%%%%%%%%%%%%%%%%%%%%%%%%%%%%%%%%%%%%%%%%%%%%%%%%%%%%
\begin{frame}[fragile]\frametitle{}

\begin{center}
{\Large Basic Demo (1 Hour)}
\end{center}
\end{frame}

%%%%%%%%%%%%%%%%%%%%%%%%%%%%%%%%%%%%%%%%%%%%%%%%%%%%%%%%%%%
%% WORKSHOP OVERVIEW
%%%%%%%%%%%%%%%%%%%%%%%%%%%%%%%%%%%%%%%%%%%%%%%%%%%%%%%%%%%
\begin{frame}[fragile]\frametitle{Workshop Overview}
\textbf{What we will build:} A chatbot with two interfaces: CLI and Streamlit, powered by LangChain + Groq LLM. We'll deploy it live to Vercel.

\textbf{What we will learn:}
  \begin{itemize}
    \item How to use Claude Code across the \textbf{full SDLC}: from specs to deployment
    \item \textbf{CLAUDE.md}, \textbf{custom commands}, \textbf{skills}, \textbf{subagents}, \textbf{hooks}, \textbf{MCP servers}, \textbf{plugins}
    \item Writing code, tests, and docs: \textit{all from the terminal}
    \item How Plan mode, permission system, and hooks enforce quality automatically
    \item \textbf{Deploying to production} using GitHub and Vercel integration
  \end{itemize}

\textbf{Duration:} Basic (Stages 1--4): \textasciitilde{}60 min. \quad Advanced (Stages 5--14): 90+ min.
\end{frame}


%%%%%%%%%%%%%%%%%%%%%%%%%%%%%%%%%%%%%%%%%%%%%%%%%%%%%%%%%%%%%%%%%%%%%%%%%%%%%%%%%%
\begin{frame}[fragile]\frametitle{SDLC Stages We Will Cover}
\begin{center}
\adjustbox{max width=\linewidth}{%
\begin{tabular}{|p{0.7cm}|l|l|}
\hline
\textbf{Stage} & \textbf{Activity} & \textbf{Claude Code Feature} \\
\hline
1 & Project bootstrap \& init & \texttt{/init}, CLAUDE.md \\
2 & Requirements \& specs & Custom subagent, Markdown \\
3 & Architecture design & Plan mode, \texttt{@workspace} \\
4 & Implementation & Auto-accept, \texttt{@file:} \\
5 & Testing (unit + synthetic) & Custom command, MCP \\
6 & Debugging & Debug subagent, context tools \\
7 & Refactoring & Skills, Rules \\
8 & Documentation & Docs subagent \\
9 & Code review & Review subagent \\
10 & GitHub workflow & MCP GitHub server \\
11 & Hooks \& automation & Hooks system (new!) \\
12 & .claude/ production setup & Rules, deny perms, skills \\
13 & Plugins & Plugin packaging (new!) \\
14 & Deployment to Vercel & GitHub integration, live hosting \\
\hline
\end{tabular}
}
\end{center}

{\small \textbf{Basic} (this section): Stages 1--4. \quad \textbf{Advanced} (next section): Stages 5--14. Each stage = one demo block.}
\end{frame}


%%%%%%%%%%%%%%%%%%%%%%%%%%%%%%%%%%%%%%%%%%%%%%%%%%%%%%%%%%%
%% STAGE 0: SETUP
%%%%%%%%%%%%%%%%%%%%%%%%%%%%%%%%%%%%%%%%%%%%%%%%%%%%%%%%%%%
\begin{frame}[fragile]\frametitle{}
\begin{center}
{\Large Demo Setup}
\end{center}
\end{frame}

%%%%%%%%%%%%%%%%%%%%%%%%%%%%%%%%%%%%%%%%%%%%%%%%%%%%%%%%%%%%%%%%%%%%%%%%%%%%%%%%%%
\begin{frame}[fragile]\frametitle{Pre-requisites Checklist}
Before starting, verify all of the following for the \textbf{mychatbot} demo:

\textbf{Essential Software} (install from Getting Started section if missing):
  \begin{itemize}
    \item Start with empty folder, switch to that directory
    \item Conda env installed, activated: \lstinline|conda --version|
    \item Claude Code installed: \lstinline|claude doctor|
    \item Node.js $\geq$ 18: \lstinline|node --version|
    \item Git configured: \lstinline|git config --global user.name "Your Name"|
    \item Python $\geq$ 3.10: \lstinline|python --version|
  \end{itemize}

\textbf{Demo-Specific Accounts \& Keys:}
  \begin{itemize}
    \item \textbf{Claude Pro/Max or API key} (already done in Getting Started)
    \item \textbf{GitHub account} and \textbf{Vercel account} (same email)
    \item Groq API key: \href{https://console.groq.com}{\texttt{console.groq.com}}
    \item GitHub Personal Access Token for MCP integration
  \end{itemize}

\textbf{Set environment variables and install demo deps:}
\begin{lstlisting}
GROQ_API_KEY="gsk_..."
GITHUB_TOKEN="ghp_..."

pip install langchain langchain-groq streamlit
pip install pytest pytest-mock faker rich
\end{lstlisting}
\end{frame}


%%%%%%%%%%%%%%%%%%%%%%%%%%%%%%%%%%%%%%%%%%%%%%%%%%%%%%%%%%%%%%%%%%%%%%%%%%%%%%%%%%
\begin{frame}[fragile]\frametitle{Choosing the Right Model}
Claude Code ships three model tiers. Use the right model for the task:

\begin{center}
\adjustbox{max width=\linewidth}{%
\begin{tabular}{|l|l|l|l|}
\hline
\textbf{Model} & \textbf{Speed} & \textbf{Cost} & \textbf{Best For} \\
\hline
claude-haiku-4-5 & Fast & Cheap & Subagents, exploration, quick edits \\
claude-sonnet-4-6 & Balanced & Medium & Default: daily coding, refactors \\
claude-opus-4-6 & Thorough & Higher & Architecture, complex reasoning \\
\hline
\end{tabular}
}
\end{center}

\textbf{Switch model inside a session:}
\begin{lstlisting}
/model claude-opus-4-6

# Or launch with specific model:
claude --model claude-opus-4-6

# Or set in .claude/settings.json:
{ "model": "claude-sonnet-4-6" }
\end{lstlisting}

% \textbf{Practical guide:}
  % \begin{itemize}
    % \item \textbf{Haiku}: subagents doing file exploration, simple search, fast questions
    % \item \textbf{Sonnet}: most coding tasks, refactoring, test generation
    % \item \textbf{Opus}: architectural design, security review, complex multi-file refactors
  % \end{itemize}
\end{frame}


%%%%%%%%%%%%%%%%%%%%%%%%%%%%%%%%%%%%%%%%%%%%%%%%%%%%%%%%%%%%%%%%%%%%%%%%%%%%%%%%%%
\begin{frame}[fragile]\frametitle{IDE Integration}
\textbf{What IDE integration adds:}
  \begin{itemize}
	\item Claude Code integrates into VS Code and JetBrains IDEs.
	\item All CLAUDE.md, settings, and MCP servers are shared:
    \item Inline diff view as Claude edits files
    \item One-click approve/reject per file change
    \item Editor context automatically sent (open file, selection)
    \item Same session continues from terminal -- no context loss
  \end{itemize}

\begin{lstlisting}
## VS Code extension:

# Install from VS Code Marketplace:
# Search: "Claude Code"

# Or via terminal:
code --install-extension anthropic.claude-code

# Use inside VS Code:
# Ctrl+Shift+P -> "Claude Code: Open"
# Or: click the Claude icon in the sidebar

## JetBrains plugin:
# Install from JetBrains Marketplace:
# Settings -> Plugins -> Search "Claude Code"
# Supports: IntelliJ IDEA, PyCharm, WebStorm, etc.
\end{lstlisting}


\end{frame}

%%%%%%%%%%%%%%%%%%%%%%%%%%%%%%%%%%%%%%%%%%%%%%%%%%%%%%%%%%%
%% PROJECT BOOTSTRAP
%%%%%%%%%%%%%%%%%%%%%%%%%%%%%%%%%%%%%%%%%%%%%%%%%%%%%%%%%%%
\begin{frame}[fragile]\frametitle{}
\begin{center}
{\Large Project Bootstrap}

% {\normalsize Teaching Claude Code about your project with CLAUDE.md}
\end{center}
\end{frame}


%%%%%%%%%%%%%%%%%%%%%%%%%%%%%%%%%%%%%%%%%%%%%%%%%%%%%%%%%%%%%%%%%%%%%%%%%%%%%%%%%%
\begin{frame}[fragile]\frametitle{Run /init: Generate CLAUDE.md}
The first command in any new project. Scans the directory and writes Claude's instruction manual.

Make sure \textbf{Plan mode} is active first (press \texttt{Shift+Tab}). Inside TUI (project is still empty, .git is there though):

\textbf{Key difference from OpenCode:} In Claude Code the file is called \texttt{CLAUDE.md} (not AGENTS.md), and Claude also supports \textbf{nested} CLAUDE.md files in subdirectories for directory-specific instructions.

Claude will ask clarifying questions about the project. Answer them. Once the plan looks right, press \texttt{Shift+Tab} to switch to normal mode and approve the file creation.


\begin{lstlisting}
/init

## What /init generates:
CLAUDE.md # Project intelligence file
.claude/settings.json # Permissions, model config

## Commit immediately:
git add CLAUDE.md
git commit -m "chore: add CLAUDE.md"
\end{lstlisting}
\end{frame}


%%%%%%%%%%%%%%%%%%%%%%%%%%%%%%%%%%%%%%%%%%%%%%%%%%%%%%%%%%%%%%%%%%%%%%%%%%%%%%%%%%
\begin{frame}[fragile]\frametitle{CLAUDE.md: The Project Memory File}
CLAUDE.md is loaded automatically at the start of \textit{every} session. Inspect and extend it with project-specific rules. Three scopes exist:

\begin{lstlisting}
~/.claude/CLAUDE.md # Global: all your projects
./CLAUDE.md   # Project: everyone on the team
./.claude/settings.local.json # Local: your overrides (gitignore)

## Edit the project CLAUDE.md:

## Project: mychatbot

## Tech Stack
- Python 3.11, LangChain 0.3+, langchain-groq
- Streamlit 1.35+, Rich (CLI output), pytest + pytest-mock

## Build Commands
- Install: `pip install -e ".[dev]"`
- Tests:   `pytest tests/ -v`
- Lint:    `ruff check src/`
- Format:  `black src/`

## Conventions
- Type hints on all functions
- Docstrings (Google style) on all public functions
- No hardcoded API keys: always read from env vars
- Never catch bare except

## File Map
- src/mychatbot/chain.py   : LangChain chain factory
- src/mychatbot/cli.py     : Rich CLI chatbot
- src/mychatbot/app.py     : Streamlit web chatbot
- tests/                   : pytest test suite
\end{lstlisting}
\end{frame}


%%%%%%%%%%%%%%%%%%%%%%%%%%%%%%%%%%%%%%%%%%%%%%%%%%%%%%%%%%%%%%%%%%%%%%%%%%%%%%%%%%
\begin{frame}[fragile]\frametitle{Auto Memory: Claude Learns As It Works}
\textbf{A unique Claude Code feature:} Claude automatically builds memory as it works, saving learnings across sessions without you writing anything.

\begin{lstlisting}
# During a session, if Claude discovers something useful: Claude writes this to CLAUDE.md automatically.
# You can also instruct Claude to remember things: "Remember that we use Black for formatting, not autopep8. Add this to CLAUDE.md."

# View what Claude has learned:
/memory

# Clear auto-generated memories:
/memory clear


## Rules files (alternative to CLAUDE.md for team rules):

# Create targeted rule files:
mkdir -p .claude/rules
cat > .claude/rules/python-style.md << 'EOF'
paths: ["src/**/*.py", "tests/**/*.py"]
Apply these rules to ALL Python files:
- Use Black formatting (line length: 88)
- Prefer dataclasses over TypedDicts for mutable data
- Always add __all__ to module-level exports
EOF
\end{lstlisting}
\end{frame}


%%%%%%%%%%%%%%%%%%%%%%%%%%%%%%%%%%%%%%%%%%%%%%%%%%%%%%%%%%%
%% REQUIREMENTS AND SPECS
%%%%%%%%%%%%%%%%%%%%%%%%%%%%%%%%%%%%%%%%%%%%%%%%%%%%%%%%%%%
\begin{frame}[fragile]\frametitle{}
\begin{center}
{\Large Requirements \& Specification}

% {\normalsize Using a custom subagent to write elaborate specs}
\end{center}
\end{frame}

% Customization overview omitted -- covered in detail in the Introduction section.

%%%%%%%%%%%%%%%%%%%%%%%%%%%%%%%%%%%%%%%%%%%%%%%%%%%%%%%%%%%%%%%%%%%%%%%%%%%%%%%%%%
\begin{frame}[fragile]\frametitle{Create a Specs Subagent}
Main Agent is Claude Code itself with its own context, tools etc. All other workers need to be SubAgents.

In Claude Code, subagents live in \texttt{.claude/agents/}. Create a dedicated \textbf{Specs subagent}:

\textbf{Key Claude Code subagent features:}
  \begin{itemize}
    \item Subagents run in \textbf{isolated contexts}: they don't bloat your main conversation
    \item \texttt{allowed-tools} restricts what the subagent can do (safety boundary)
    \item Claude automatically selects a subagent when its \texttt{description} matches the task
  \end{itemize}

\begin{lstlisting}
mkdir -p .claude/agents

cat > .claude/agents/specs.md << 'EOF'
---
name: specs
description: >
  Technical specification writer. Produces clear,
  structured Markdown spec docs from user stories or
  feature descriptions. Read-only: never modifies source
  code. Use when writing a technical spec for any feature.
model: claude-opus-4-6
allowed-tools: Read, Write, Glob
---
You are a senior technical writer and software architect.
When asked to write a spec, produce a Markdown document
covering: Overview, Goals, Non-Goals, User Stories,
Functional Requirements, Non-Functional Requirements,
Data Models, API/Interface Contracts, and Open Questions.
Always save output to docs/specs/<feature>.md.
EOF
\end{lstlisting}


\end{frame}


%%%%%%%%%%%%%%%%%%%%%%%%%%%%%%%%%%%%%%%%%%%%%%%%%%%%%%%%%%%%%%%%%%%%%%%%%%%%%%%%%%
\begin{frame}[fragile]\frametitle{Generate the Project Spec}
Inside the interactive REPL (Read-Eval-Print Loop, text box where you can type and execute code, for now, here its the normal mode ie confirm before accept), invoke the specs subagent:

\begin{lstlisting}
Write a technical specification for mychatbot.
It should cover:
- The CLI chatbot interface (multi-turn conversation,
  history, exit command)
- The Streamlit web chatbot (session state, streaming)
- The shared LangChain chain module (model config,
  system prompt injection, memory management)
- Environment variable requirements
- Error handling expectations
Save to docs/specs/mychatbot-spec.md

## Claude detects this matches the ``specs'' subagent description and routes automatically. You can also invoke explicitly:


# Explicit subagent invocation:
Use the specs subagent to write a spec for mychatbot.

# In non-interactive mode:
claude -p "Use the specs subagent to write a spec for
  mychatbot. Save to docs/specs/mychatbot-spec.md"

## Git
git add docs/
git commit -m "docs: add mychatbot spec"
\end{lstlisting}
\end{frame}


%%%%%%%%%%%%%%%%%%%%%%%%%%%%%%%%%%%%%%%%%%%%%%%%%%%%%%%%%%%
%% ARCHITECTURE
%%%%%%%%%%%%%%%%%%%%%%%%%%%%%%%%%%%%%%%%%%%%%%%%%%%%%%%%%%%
\begin{frame}[fragile]\frametitle{}
\begin{center}
{\Large Architecture Design}

% {\normalsize Plan mode: read-only exploration before any edits}
\end{center}
\end{frame}


%%%%%%%%%%%%%%%%%%%%%%%%%%%%%%%%%%%%%%%%%%%%%%%%%%%%%%%%%%%%%%%%%%%%%%%%%%%%%%%%%%
\begin{frame}[fragile]\frametitle{Plan Mode: Design Before Coding}
Plan mode restricts Claude to \textbf{read-only exploration} -- no file edits, no bash execution:

\textbf{After plan review, three approval options:}
  \begin{itemize}
    \item ``Yes, clear context and auto-accept edits'' (recommended: fresh context for implementation)
    \item ``Yes, and manually approve edits'' (review each change)
    \item ``Yes, auto-accept edits'' (run with preserved context)
  \end{itemize}

\begin{lstlisting}
# Enter plan mode (cycles through modes):
Shift+Tab # Cycles: normal -> plan -> auto-accept

# Or via command:
/plan          # Enter plan mode
/plan refactor the auth module  # With description

## In Plan mode, ask for the architecture:

Read docs/specs/mychatbot-spec.md, then design the
project structure for mychatbot:
- Package layout under src/mychatbot/
- Module responsibilities (chain.py, cli.py, app.py, config.py)
- Data flow between modules
- Dependency graph
- What should be tested vs. what should be mocked
Show me the plan before creating any files.
\end{lstlisting}


\end{frame}


%%%%%%%%%%%%%%%%%%%%%%%%%%%%%%%%%%%%%%%%%%%%%%%%%%%%%%%%%%%%%%%%%%%%%%%%%%%%%%%%%%
\begin{frame}[fragile]\frametitle{Plan Mode: Extended Thinking}
For complex architectural decisions, combine Plan mode with \textbf{Extended Thinking}.

\textbf{Pro Tip:} Extended thinking is especially valuable in Plan mode because it can explore multiple approaches without consuming execution tokens on wasted code.

\begin{lstlisting}
# Toggle extended thinking inside a session:
Alt+T  # Toggle extended thinking on/off

# Or launch with thinking:
claude "Design the architecture for mychatbot"
# Then press Alt+T before sending complex questions

## In Plan mode with extended thinking enabled:
Given docs/specs/mychatbot-spec.md, should we use:
1. LangChain ConversationBufferMemory (simple, risk of token bloat), or
2. LangChain ConversationSummaryMemory (summarises old turns, lower token use)?

Consider: multi-turn history size, streaming support, Groq API limits, testability. Think carefully and give a recommendation with trade-offs.
\end{lstlisting}

\end{frame}


%%%%%%%%%%%%%%%%%%%%%%%%%%%%%%%%%%%%%%%%%%%%%%%%%%%%%%%%%%%%%%%%%%%%%%%%%%%%%%%%%%
\begin{frame}[fragile]\frametitle{Scaffold the Project Structure}
Exit plan mode (press \texttt{Shift+Tab}). In normal mode:

\begin{lstlisting}
Create the mychatbot package scaffold based on the architecture plan and add stub methods:

## Result may be:
- src/mychatbot/__init__.py
- src/mychatbot/chain.py     (stub: make_chain())
- src/mychatbot/cli.py       (stub: main())
- src/mychatbot/app.py       (stub: Streamlit skeleton)
- src/mychatbot/config.py    (stub: Config dataclass)
- tests/__init__.py
- tests/test_chain.py        (stub)
- pyproject.toml             (project metadata + deps)
- README.md                  (one-line placeholder)

## Git ops:
git add .
git commit -m "chore: scaffold mychatbot package"

## Context management tip:
/cost       # Check token usage so far
/compact    # Summarise older context to free space
\end{lstlisting}
\end{frame}


%%%%%%%%%%%%%%%%%%%%%%%%%%%%%%%%%%%%%%%%%%%%%%%%%%%%%%%%%%%
%% IMPLEMENTATION
%%%%%%%%%%%%%%%%%%%%%%%%%%%%%%%%%%%%%%%%%%%%%%%%%%%%%%%%%%%
\begin{frame}[fragile]\frametitle{}
\begin{center}
{\Large Implementation}

{\normalsize File references, auto-accept mode, and session management}
\end{center}
\end{frame}


%%%%%%%%%%%%%%%%%%%%%%%%%%%%%%%%%%%%%%%%%%%%%%%%%%%%%%%%%%%%%%%%%%%%%%%%%%%%%%%%%%
\begin{frame}[fragile]\frametitle{Review and Implement chain.py}
In normal mode (use \texttt{@file:} to reference files precisely). Review the contents first and accordingly ask to add the remaining.

\begin{lstlisting}
@file:src/mychatbot/chain.py
@file:docs/specs/mychatbot-spec.md
Implement chain.py:
- make_chain(model: str, system_prompt: str) -> Callable
- Use ChatGroq from langchain_groq
- Accept conversation history as list[dict] (role/content)
- Return a callable that takes history and returns str
- Read GROQ_API_KEY from env via Config
- Raise EnvironmentError if key is missing
- Include full type hints and Google-style docstrings
- Handle Groq API errors with informative messages


## Reviewing Claude's choices:

# Ask Claude to explain its implementation decisions:
Why did you use ChatPromptTemplate instead of
MessagesPlaceholder? Is there a better approach
for multi-turn history?

## Git Ops
git add src/mychatbot/chain.py
git commit -m "feat: implement LangChain chain factory"
\end{lstlisting}
\end{frame}


%%%%%%%%%%%%%%%%%%%%%%%%%%%%%%%%%%%%%%%%%%%%%%%%%%%%%%%%%%%%%%%%%%%%%%%%%%%%%%%%%%
\begin{frame}[fragile]\frametitle{Review and Implement cli.py ,  app.py}
Review the contents first and accordingly ask to add the remaining.


\begin{lstlisting}
## CLI implementation:

@file:src/mychatbot/cli.py
Implement the CLI chatbot:
- main() as entry point
- Use Rich: Console, Panel (welcome), Markdown (responses)
- [You] prompt styled, [mychatbot] in green
- Maintain conversation history as list[dict]
- Commands: /exit or /quit to end, /clear to reset history
- Ctrl+C gracefully exits with "Goodbye!" message
- Call make_chain() from mychatbot.chain
- Model/system prompt configurable via env vars


## Streamlit implementation:

@file:src/mychatbot/app.py
Implement the Streamlit chatbot:
- st.set_page_config: title="mychatbot"
- st.session_state["messages"] for conversation history
- st.session_state["chain"] to cache the chain callable
- st.chat_message bubbles for conversation display
- Sidebar: model selector, system prompt text area
- Spinner while waiting for the model response
- Handle missing GROQ_API_KEY with st.error and st.stop()
\end{lstlisting}
\end{frame}


%%%%%%%%%%%%%%%%%%%%%%%%%%%%%%%%%%%%%%%%%%%%%%%%%%%%%%%%%%%%%%%%%%%%%%%%%%%%%%%%%%
\begin{frame}[fragile]\frametitle{Session Management}

\begin{lstlisting}
# Resume most recent session after closing terminal:
claude -c

# Name the current session for later:
/rename mychatbot-implementation

# Resume a named session:
claude --resume mychatbot-implementation

# Fork a session to try an alternative approach:
claude --resume mychatbot-implementation --fork-session

## Run the chatbot to verify:
python -m mychatbot.cli     # CLI interface
streamlit run src/mychatbot/app.py  # http://localhost:8501
\end{lstlisting}
\end{frame}
